\section{Kiến trúc Transformer}
\label{sec:transformer}

\subsection{Giới thiệu}

Transformer là kiến trúc mạng nơ-ron được giới thiệu bởi Vaswani et al. năm 2017 trong bài báo "Attention is All You Need" \cite{vaswani2017attention}. Khác với RNN và LSTM xử lý chuỗi tuần tự, Transformer sử dụng cơ chế self-attention để xử lý toàn bộ chuỗi song song, cho phép nắm bắt long-range dependencies hiệu quả hơn.

\subsection{Self-Attention Mechanism}

\subsubsection{Scaled Dot-Product Attention}

Cho một chuỗi đầu vào, self-attention tính toán representation của mỗi vị trí dựa trên toàn bộ chuỗi:

\begin{equation}
\text{Attention}(\mathbf{Q}, \mathbf{K}, \mathbf{V}) = \text{softmax}\left(\frac{\mathbf{Q}\mathbf{K}^T}{\sqrt{d_k}}\right)\mathbf{V}
\end{equation}

trong đó:
\begin{itemize}
    \item $\mathbf{Q} = \mathbf{X}\mathbf{W}^Q$ (Query): "Tôi đang tìm kiếm thông tin gì?"
    \item $\mathbf{K} = \mathbf{X}\mathbf{W}^K$ (Key): "Tôi chứa thông tin gì?"
    \item $\mathbf{V} = \mathbf{X}\mathbf{W}^K$ (Value): "Thông tin thực tế tôi cung cấp"
    \item $d_k$: Dimensionality của key vectors (scaling factor)
\end{itemize}

Attention weights $\alpha_{ij} = \text{softmax}\left(\frac{\mathbf{q}_i \mathbf{k}_j^T}{\sqrt{d_k}}\right)$ đo lường mức độ liên quan giữa vị trí $i$ và $j$.

\subsubsection{Multi-Head Attention}

Multi-head attention cho phép mô hình chú ý đến các khía cạnh khác nhau của thông tin:

\begin{align}
\text{MultiHead}(\mathbf{Q}, \mathbf{K}, \mathbf{V}) &= \text{Concat}(\text{head}_1, \ldots, \text{head}_h)\mathbf{W}^O \\
\text{head}_i &= \text{Attention}(\mathbf{Q}\mathbf{W}_i^Q, \mathbf{K}\mathbf{W}_i^K, \mathbf{V}\mathbf{W}_i^V)
\end{align}

với $h$ là số attention heads. DNABERT-2 sử dụng 12 attention heads.

\subsection{Positional Encoding}

Vì self-attention không có khái niệm về thứ tự, cần thêm positional information:

\subsubsection{Sinusoidal Positional Encoding}

Transformer gốc sử dụng sinusoidal functions:

\begin{align}
PE_{(pos, 2i)} &= \sin\left(\frac{pos}{10000^{2i/d}}\right) \\
PE_{(pos, 2i+1)} &= \cos\left(\frac{pos}{10000^{2i/d}}\right)
\end{align}

\subsubsection{Learned Positional Embeddings}

BERT học positional embeddings như parameters, nhưng bị giới hạn bởi maximum sequence length.

\subsubsection{Attention with Linear Biases (ALiBi)}

DNABERT-2 sử dụng ALiBi \cite{press2021train}, thêm bias vào attention scores:

\begin{equation}
\text{softmax}\left(\mathbf{q}_i \mathbf{K}^T + m \cdot [-(i-1), \ldots, -1, 0]\right)
\end{equation}

với $m$ là slope khác nhau cho mỗi head. ALiBi cho phép xử lý chuỗi có độ dài tùy ý mà không cần retrain.

\subsection{Feed-Forward Networks}

Mỗi transformer layer có position-wise feed-forward network:

\begin{equation}
\text{FFN}(\mathbf{x}) = \max(0, \mathbf{x}\mathbf{W}_1 + \mathbf{b}_1)\mathbf{W}_2 + \mathbf{b}_2
\end{equation}

hoặc với GELU activation:

\begin{equation}
\text{FFN}(\mathbf{x}) = \text{GELU}(\mathbf{x}\mathbf{W}_1 + \mathbf{b}_1)\mathbf{W}_2 + \mathbf{b}_2
\end{equation}

\subsection{Transformer Encoder}

Một transformer encoder layer bao gồm:

\begin{align}
\mathbf{h}' &= \text{LayerNorm}(\mathbf{x} + \text{MultiHeadAttention}(\mathbf{x})) \\
\mathbf{h} &= \text{LayerNorm}(\mathbf{h}' + \text{FFN}(\mathbf{h}'))
\end{align}

với residual connections và layer normalization.

\subsection{Ưu điểm của Transformer}

\begin{itemize}
    \item \textbf{Parallelization}: Xử lý toàn bộ chuỗi song song, nhanh hơn RNN
    \item \textbf{Long-range dependencies}: Self-attention trực tiếp kết nối mọi cặp vị trí
    \item \textbf{Flexible attention patterns}: Mỗi head có thể học các patterns khác nhau
    \item \textbf{Scalability}: Hiệu suất cải thiện với model size và data
\end{itemize}

\subsection{Transformer cho genomics}

Transformer đã được áp dụng thành công cho nhiều bài toán genomics:
\begin{itemize}
    \item DNABERT, DNABERT-2: Pre-trained models cho DNA sequences
    \item Nucleotide Transformer: Large-scale genome foundation model
    \item ProtTrans: Protein language models
    \item RNA-FM: RNA sequence modeling
\end{itemize}

Khả năng nắm bắt long-range dependencies của transformer đặc biệt phù hợp với dữ liệu genome, nơi các yếu tố điều hòa có thể cách xa nhau hàng nghìn base pairs.
